\section{Introduction}

The goal of the \emph{Rethinking Spreadsheets} project is to reconsider the
interface that existing spreadsheets provide to the user, and to improve upon
it.

The traditional way to program a spreadsheet is to type a formula into a cell;
the formula represents a calculation based on the values in other cells.  The
cell then displays the value of the formula.  Often, the user copies the
formula in one cell into adjacent cells; the application adjusts references to
other cells appropriately.

The difficulty with this approach is that it is very hard to scrutinise the
code.  The user has to click on each cell individually and read the formula.
If is very easy to make mistakes.  

For example, the Economists Reinhart and Rogoff claimed\footnote{Carmen
  M. Reinhart and Kenneth S. Rogoff, "Growth in a Time of Debt". American
  Economic Review. 100 (2): 573–78, 2010.} that if a country has national debt
of over 90\% of GDP, it leads to much reduced growth.  Various nations adopted
policies of austerity based on this paper.  However, Herdon, Ash and
Pollin\footnote{Thomas Herndon, Michael Ash, Robert Pollin. "Does High Public
  Debt Consistently Stifle Economic Growth? A Critique of Reinhart and
  Rogoff". Political Economy Research Institute, University of Massachusetts
  Amherst.  2013.} found that the result was incorrect, partly because of an
error in the spreadsheet which omitted some countries.  Such errors are easy
to make, but hard to spot.

Our thesis is that there should be a separation between the rules that
define calculations and the results of those calculations.  The rules are
given in a separate script; and the results appear in the spreadsheet.

In order to allow scrutiny of the script, and to support maintenance, it
should be a reasonable length.  If similar calculations are performed to
produce multiple cells, the relevant formula should be defined once, and then
applied to relevant cells via an iteration: it should not be necessary to copy
the formula. 

Note, though, that we propose keeping input data and calculated cells as part
of the same spreadsheet: it is often convenient to view those together.

Further, the language used for most existing spreadsheets---simple functions
that are combined together---is not convenient for anything other than very
simple applications.  Spreadsheet users should be able to use a language
similar to modern programming languages. 
%
The script language we propose is a declarative language, influenced by both
Haskell and the functional fragment of Scala.  In particular, the language
allows type checking of the script.

A script contains a combination of declarations (of values and functions), and
commands (e.g.~to write a value into a cell). 

In Section~\ref{sec:example} we present a simple example script, to introduce
the main parts of the language.  In Section~\ref{sec:using} we describe how
to use the spreadsheet application.   Then in Section~\ref{sec:language} we
present a more detailed description of the language.  We assume some
familiarity with programming. 
